\documentclass[11pt,a4paper]{article}
\usepackage[utf8]{inputenc}
\usepackage[UTF8]{ctex} % 支持中文显示
\usepackage{amsmath, amssymb, amsfonts}
\usepackage{geometry}
\usepackage{bm}
\usepackage{tcolorbox}
\usepackage[
    colorlinks=true,
    linkcolor=blue,
    unicode=true,          % 关键：使用 unicode 编码处理书签
    pdfpagemode=UseNone,   % 避免打开时自动弹出书签栏导致某些预览器崩溃
    pdfstartview=FitH
]{hyperref}
\geometry{left=2cm, right=2cm, top=2.5cm, bottom=2.5cm}

\title{分析力学笔记整理 (Analytical Mechanics)}
\author{魏舟}
\date{\today}

\begin{document}

\maketitle
\newpage
\tableofcontents
\newpage
\section{广义坐标与约束 (Generalized Coordinates \& Constraints)}

对于具有 $N$ 个质点的系统，其位置由 $3N$ 个直角坐标确定。引入广义坐标：
\begin{align*}
    q_1 &= q_1(x_1, y_1, z_1, x_2, y_2, \dots, z_N, t) \\
    q_2 &= q_2(x_1, y_1, z_1, \dots, z_N, t) \\
    &\vdots \\
    q_{3N} &= q_{3N}(x_1, y_1, z_1, \dots, z_N, t)
\end{align*}

若用 $(x_1, y_1, z_1, \dots, z_N)$ 间 $(q_1, q_2, \dots, q_n)$ 表示，则 \textbf{Jacobian 行列式}为：
$$ \frac{\partial(q_1, q_2, \dots, q_n)}{\partial(x_1, x_2, \dots, x_n)} = 
\begin{vmatrix} 
\frac{\partial q_1}{\partial x_1} & \frac{\partial q_1}{\partial x_2} & \dots & \frac{\partial q_1}{\partial x_n} \\
\vdots & \vdots & \ddots & \vdots \\
\frac{\partial q_n}{\partial x_1} & \frac{\partial q_n}{\partial x_2} & \dots & \frac{\partial q_n}{\partial x_n} 
\end{vmatrix} \neq 0 $$
这意味着广义坐标之间是\textbf{线性无关 (Linearly independent)}的。

\section*{广义坐标独立性的数学推导}

在分析力学或多变量微积分中，若一组变量 $q_1, q_2, \dots, q_n$ 关于 $x_1, x_2, \dots, x_n$ 的雅可比行列式（Jacobian determinant）不为零：
\[
J = \det\left(\frac{\partial q_i}{\partial x_j}\right) = 
\begin{vmatrix} 
\frac{\partial q_1}{\partial x_1} & \frac{\partial q_1}{\partial x_2} & \dots & \frac{\partial q_1}{\partial x_n} \\
\vdots & \vdots & \ddots & \vdots \\
\frac{\partial q_n}{\partial x_1} & \frac{\partial q_n}{\partial x_2} & \dots & \frac{\partial q_n}{\partial x_n} 
\end{vmatrix} \neq 0 
\]
这意味着广义坐标之间是\textbf{相互独立（Functionally Independent）}且\textbf{线性无关}的。

\subsection*{推导过程}

我们可以利用反证法结合全微分性质来证明：

1. **假设相关性**：若广义坐标 $q_1, \dots, q_n$ 之间线性相关，则存在一组不全为零的常数 $c_1, c_2, \dots, c_n$，使得：
   \[ \sum_{i=1}^n c_i dq_i = 0 \]

2. **展开全微分**：由于 $q_i$ 是 $x_j$ 的函数，其全微分为 $dq_i = \sum_{j=1}^n \frac{\partial q_i}{\partial x_j} dx_j$。代入上式得：
   \[ \sum_{i=1}^n c_i \left( \sum_{j=1}^n \frac{\partial q_i}{\partial x_j} dx_j \right) = 0 \]

3. **变换求和次序**：
   \[ \sum_{j=1}^n \left( \sum_{i=1}^n c_i \frac{\partial q_i}{\partial x_j} \right) dx_j = 0 \]
   由于自变量 $x_j$ 是彼此独立的（即 $dx_j$ 线性无关），要使上式对任意 $dx_j$ 成立，括号内的系数必须全部为零：
   \[ \sum_{i=1}^n c_i \frac{\partial q_i}{\partial x_j} = 0, \quad (j=1, 2, \dots, n) \]

4. **矩阵表述**：上述方程组可以表示为线性方程组 $\mathbf{J}^T \mathbf{c} = \mathbf{0}$：
   \[
   \begin{pmatrix} 
   \frac{\partial q_1}{\partial x_1} & \dots & \frac{\partial q_n}{\partial x_1} \\
   \vdots & \ddots & \vdots \\
   \frac{\partial q_1}{\partial x_n} & \dots & \frac{\partial q_n}{\partial x_n} 
   \end{pmatrix}
   \begin{pmatrix} c_1 \\ \vdots \\ c_n \end{pmatrix} = \mathbf{0}
   \]
   
5. **结论**：若 $\det(\mathbf{J}) \neq 0$（即雅可比行列式不为零），则该齐次线性方程组只有唯一零解 $\mathbf{c} = \mathbf{0}$。这与“存在不全为零的常数”的假设矛盾。

因此，$\det(\mathbf{J}) \neq 0$ 是广义坐标之间\textbf{线性无关}的充分必要条件.



\subsection{约束 (Constraints)}
\begin{itemize}
    \item \textbf{自由度：} 完全指定系统每个部分的位置所需要的独立广义坐标的数目。
    \item 每个约束减少一个自由度。
    \item \textbf{完整约束 (Holonomic)：} 如果约束方程可表示为 $f(q_1, q_2, \dots, q_n, t) = 0$。
\end{itemize}

\section{速度、虚位移与虚功}

\subsection{广义速度}
设 $x_i = x_i(q_1, q_2, \dots, q_n, t)$，则速度 $v_i = \frac{dx_i}{dt}$：
$$ \frac{dx_i}{dt} = \sum_{k} \frac{\partial x_i}{\partial q_k} \dot{q}_k + \frac{\partial x_i}{\partial t} $$
由此推导出重要关系：
$$ \frac{\partial v_i}{\partial \dot{q}_j} = \frac{\partial}{\partial \dot{q}_j} \left( \sum_k \frac{\partial x_i}{\partial q_k} \dot{q}_k + \frac{\partial x_i}{\partial t} \right) = \frac{\partial x_i}{\partial q_j} $$

\subsection{虚位移与虚功原理}
虚位移 $\delta x_i$ (在 $t$ 不变的情况下)：$\delta x_i = \sum_{\alpha=1}^{n} \frac{\partial x_i}{\partial q_\alpha} \delta q_\alpha$。
虚功：
$$ \delta W = \sum_{i=1}^{3N} F_i \delta x_i = \sum_{\alpha=1}^n \left( \sum_{i=1}^{3N} F_i \frac{\partial x_i}{\partial q_\alpha} \right) \delta q_\alpha $$
定义\textbf{广义力}为：
\begin{tcolorbox}[colframe=red!50!black, colback=white, title=定义]
    $Q_\alpha = \sum_{i=1}^{3N} F_i \frac{\partial x_i}{\partial q_\alpha}$
\end{tcolorbox}
则虚功方程为：$\delta W = \sum_{\alpha=1}^n Q_\alpha \delta q_\alpha$。平衡时 $\delta W = 0$。

\section{拉格朗日动力学 (Lagrangian Dynamics)}

定义拉格朗日量为动能与势能之差：$\boxed{L = T - V}$。
\textbf{Lagrange 方程：}
$$ \boxed{\frac{d}{dt} \left( \frac{\partial L}{\partial \dot{q}_i} \right) - \frac{\partial L}{\partial q_i} = 0} $$

\subsection{实例：阿特伍德机 (Atwood's machine)}

设绳长为 $L$，$x_1 + x_2 = L$。
$T = \frac{1}{2}(m_1 + m_2)\dot{x}_1^2$，$V = -m_1 g x_1 - m_2 g (L - x_1)$。
$L = \frac{1}{2}(m_1+m_2)\dot{x}_1^2 + (m_1-m_2)g x_1 + m_2 g L$。
由拉格朗日方程：
$$ (m_1+m_2)\ddot{x}_1 - (m_1-m_2)g = 0 \implies \ddot{x}_1 = \frac{m_1-m_2}{m_1+m_2}g $$


\section{守恒定律 (Conservation Laws) 及其推导}

\subsection{空间均匀性与动量守恒 (Homogeneity of Space)}
若物理规律在空间平移下不变，即拉格朗日量 $L$ 在坐标微移 $\delta \vec{r}_i = \vec{\epsilon}$（$\vec{\epsilon}$ 为任意常矢量）下保持不变，则：
\[ \delta L = \sum_i \frac{\partial L}{\partial \vec{r}_i} \cdot \delta \vec{r}_i = \vec{\epsilon} \cdot \sum_i \frac{\partial L}{\partial \vec{r}_i} = 0 \]
由于 $\vec{\epsilon}$ 是任意的，必有 $\sum_i \frac{\partial L}{\partial \vec{r}_i} = 0$。根据拉格朗日方程：
\[ \frac{d}{dt} \left( \frac{\partial L}{\partial \dot{\vec{r}}_i} \right) = \frac{\partial L}{\partial \vec{r}_i} \]
代入上式得：
\[ \frac{d}{dt} \sum_i \frac{\partial L}{\partial \dot{\vec{r}}_i} = \frac{d}{dt} \sum_i \vec{p}_i = 0 \]
由此得出总动量 $\vec{P} = \sum \vec{p}_i$ 为常矢量。即：\textbf{空间均匀性对应动量守恒}。

\subsection{空间各向同性与角动量守恒 (Isotropy of Space)}
若物理规律在绕任意轴的无穷小转动 $\delta \vec{r}_i = \delta \vec{\phi} \times \vec{r}_i$ 下不变，则：
\[ \delta L = \sum_i \left( \frac{\partial L}{\partial \vec{r}_i} \cdot \delta \vec{r}_i + \frac{\partial L}{\partial \dot{\vec{r}}_i} \cdot \delta \dot{\vec{r}}_i \right) = 0 \]
利用 $\delta \dot{\vec{r}}_i = \delta \vec{\phi} \times \dot{\vec{r}}_i$ 以及 $\dot{\vec{p}}_i = \frac{\partial L}{\partial \vec{r}_i}$，代入上式：
\[ \sum_i \left( \dot{\vec{p}}_i \cdot (\delta \vec{\phi} \times \vec{r}_i) + \vec{p}_i \cdot (\delta \vec{\phi} \times \dot{\vec{r}}_i) \right) = 0 \]
利用矢量混合积的循环对称性 $\vec{A} \cdot (\vec{B} \times \vec{C}) = \vec{B} \cdot (\vec{C} \times \vec{A})$：
\[ \delta \vec{\phi} \cdot \sum_i \left( \vec{r}_i \times \dot{\vec{p}}_i + \dot{\vec{r}}_i \times \vec{p}_i \right) = \delta \vec{\phi} \cdot \frac{d}{dt} \left( \sum_i \vec{r}_i \times \vec{p}_i \right) = 0 \]
由于 $\delta \vec{\phi}$ 的任意性，得出总角动量 $\vec{L}_{tot} = \sum \vec{r}_i \times \vec{p}_i$ 守恒。即：\textbf{空间各向同性对应角动量守恒}。

\subsection{时间均匀性与能量守恒 (Homogeneity of Time)}
若拉格朗日量不显含时间（$\frac{\partial L}{\partial t} = 0$），考虑 $L(q_i, \dot{q}_i)$ 对时间的总导数：
\[ \frac{dL}{dt} = \sum_i \frac{\partial L}{\partial q_i} \dot{q}_i + \sum_i \frac{\partial L}{\partial \dot{q}_i} \ddot{q}_i \]
利用拉格朗日方程 $\frac{\partial L}{\partial q_i} = \frac{d}{dt} \frac{\partial L}{\partial \dot{q}_i} = \dot{p}_i$：
\[ \frac{dL}{dt} = \sum_i \dot{p}_i \dot{q}_i + \sum_i p_i \ddot{q}_i = \frac{d}{dt} \left( \sum_i p_i \dot{q}_i \right) \]
整理得：
\[ \frac{d}{dt} \left( \sum_i p_i \dot{q}_i - L \right) = 0 \]
定义能量函数（哈密顿量） $h = \sum p_i \dot{q}_i - L$，则 $h$ 为守恒量。即：\textbf{时间均匀性对应能量守恒}。

\subsection{维里定理 (The Virial Theorem)}
\subsection*{维里定理 (The Virial Theorem) 的推导}

定义系统的“维里函数” (Virial) $G$ 为：
\[ G = \sum_k \vec{p}_k \cdot \vec{r}_k \]
其中 $\vec{p}_k = m_k \dot{\vec{r}}_k$ 为第 $k$ 个质点的动量。对 $G$ 关于时间 $t$ 求全导数：
\[ \frac{dG}{dt} = \sum_k \dot{\vec{p}}_k \cdot \vec{r}_k + \sum_k \vec{p}_k \cdot \dot{\vec{r}}_k \]
根据牛顿第二定律 $\dot{\vec{p}}_k = \vec{F}_k$，以及动量定义 $\vec{p}_k \cdot \dot{\vec{r}}_k = m_k \dot{\vec{r}}_k \cdot \dot{\vec{r}}_k = m_k v_k^2$，上式可化为：
\[ \frac{dG}{dt} = \sum_k \vec{F}_k \cdot \vec{r}_k + \sum_k m_k v_k^2 = \sum_k \vec{F}_k \cdot \vec{r}_k + 2T \]
其中 $T$ 为系统的总动能。

现在考虑该物理量在时间周期 $\tau$ 内的平均值：
\[ \left\langle \frac{dG}{dt} \right\rangle = \frac{1}{\tau} \int_0^\tau \frac{dG}{dt} dt = \frac{G(\tau) - G(0)}{\tau} \]
对于\textbf{有界运动}（即质点的坐标与动量均在有限范围内变化），当 $\tau \to \infty$ 时，上式右端趋于零：
\[ \lim_{\tau \to \infty} \frac{G(\tau) - G(0)}{\tau} = 0 \]
因此，我们得到时间平均值的关系：
\[ \langle 2T \rangle + \left\langle \sum_k \vec{F}_k \cdot \vec{r}_k \right\rangle = 0 \]
即：
\[ \boxed{\langle 2T \rangle = - \left\langle \sum_k \vec{F}_k \cdot \vec{r}_k \right\rangle} \]

\textbf{特殊情况（势能为齐次函数）：}
若势能满足 $V(\alpha \vec{r}) = \alpha^n V(\vec{r})$（即 $n$ 次齐次函数，如引力场 $n=-1$），由欧拉齐次函数定理可知：
\[ \langle 2T \rangle = n \langle V \rangle \]
例如在万有引力场中，$n = -1$，则有 $\langle 2T \rangle = - \langle V \rangle$。
\[ \boxed{\langle 2T \rangle = - \left\langle \sum \vec{F}_k \cdot \vec{r}_k \right\rangle} \]

\section{变分计算 (The Calculus of Variations)}

变分法的核心目标是寻找一个函数 $y(x)$，使得泛函（Functional） $I$ 取得平稳值（极值）：
\[ I = \int_{x_1}^{x_2} \Phi(x, y, y') dx \]
其中 $y' = \frac{dy}{dx}$，且边界条件固定为 $y(x_1) = y_1, y(x_2) = y_2$。

\subsection{Euler-Lagrange 方程的推导}

1. **引入变分路径**：
   设 $y(x)$ 是使泛函取得极值的真实路径，引入邻近路径 $Y(x, \epsilon) = y(x) + \epsilon \eta(x)$。
   其中 $\epsilon$ 是无穷小参数，$\eta(x)$ 是任意可微函数，且由于边界固定，必须满足 $\eta(x_1) = \eta(x_2) = 0$。

2. **泛函对参数求导**：
   将 $Y$ 代入泛函，使其成为关于 $\epsilon$ 的函数 $I(\epsilon)$。极值条件要求在 $\epsilon = 0$ 处导数为零：
   \[ \frac{dI}{d\epsilon}\bigg|_{\epsilon=0} = \int_{x_1}^{x_2} \left( \frac{\partial \Phi}{\partial y} \frac{\partial Y}{\partial \epsilon} + \frac{\partial \Phi}{\partial y'} \frac{\partial Y'}{\partial \epsilon} \right) dx = 0 \]
   注意到 $\frac{\partial Y}{\partial \epsilon} = \eta(x)$ 且 $\frac{\partial Y'}{\partial \epsilon} = \eta'(x)$，上式变为：
   \[ \int_{x_1}^{x_2} \left( \frac{\partial \Phi}{\partial y} \eta(x) + \frac{\partial \Phi}{\partial y'} \eta'(x) \right) dx = 0 \]

3. **分部积分**：
   对第二项进行分部积分：
   \[ \int_{x_1}^{x_2} \frac{\partial \Phi}{\partial y'} \eta'(x) dx = \left[ \frac{\partial \Phi}{\partial y'} \eta(x) \right]_{x_1}^{x_2} - \int_{x_1}^{x_2} \eta(x) \frac{d}{dx} \left( \frac{\partial \Phi}{\partial y'} \right) dx \]
   由于 $\eta(x_1) = \eta(x_2) = 0$，边界项消失。

4. **基本引理应用**：
   代回积分式得：
   \[ \int_{x_1}^{x_2} \eta(x) \left( \frac{\partial \Phi}{\partial y} - \frac{d}{dx} \frac{\partial \Phi}{\partial y'} \right) dx = 0 \]
   根据变分法基本引理，由于 $\eta(x)$ 的任意性，括号内的表达式必须恒为零，得到 \textbf{Euler-Lagrange 方程}：
   \[ \boxed{\frac{\partial \Phi}{\partial y} - \frac{d}{dx} \left( \frac{\partial \Phi}{\partial y'} \right) = 0} \]

\subsection{贝尔特拉米恒等式 (Beltrami Identity)}
若 $\Phi$ 不显含 $x$（即 $\frac{\partial \Phi}{\partial x} = 0$），则存在第一积分。
考虑全导数 $\frac{d\Phi}{dx} = \frac{\partial \Phi}{\partial x} + \frac{\partial \Phi}{\partial y}y' + \frac{\partial \Phi}{\partial y'}y''$。
利用 E-L 方程代换 $\frac{\partial \Phi}{\partial y}$，整理可得：
\[ \boxed{\Phi - y' \frac{\partial \Phi}{\partial y'} = C \quad (\text{常数})} \]

\subsection{实例：最速降线问题 (Brachistochrone problem)}



\textbf{问题描述}：质点在重力作用下从 $A(0,0)$ 滑到 $B(x_2, y_2)$，求耗时最短的曲线 $y(x)$。
1. **建立泛函**：
   由能量守恒 $v = \sqrt{2gy}$。时间微元 $dt = \frac{ds}{v} = \frac{\sqrt{1+y'^2}dx}{\sqrt{2gy}}$。
   总时间 $T[y] = \int_0^{x_2} \sqrt{\frac{1+y'^2}{2gy}} dx$。

2. **求解变分**：
   此处 $\Phi = \sqrt{\frac{1+y'^2}{y}}$（忽略常数项）。由于 $\Phi$ 不显含 $x$，使用贝尔特拉米恒等式：
   \[ \sqrt{\frac{1+y'^2}{y}} - y' \left( \frac{y'}{\sqrt{y(1+y'^2)}} \right) = \frac{1}{\sqrt{C}} \]
   简化为：$y(1+y'^2) = C = 2a$。

3. **参数方程解**：
   令 $y' = \cot(\phi/2)$，代入微分方程并积分，可得其轨迹为\textbf{旋轮线 (Cycloid)}：
   \[ \boxed{x = a(\phi - \sin\phi), \quad y = a(1 - \cos\phi)} \]


\section{达朗贝尔原理与拉格朗日方程的导出}

\subsection{达朗贝尔原理 (d'Alembert's Principle)}
对于一个由 $N$ 个质点组成的系统，达朗贝尔原理是牛顿第二定律的一种重新表述。其核心思想是将动力学问题转化为静力学平衡问题。其数学形式为：
\[ \sum_{\alpha=1}^{N}(\vec{F}_{\alpha}^{ext}-\dot{\vec{p}}_{\alpha})\cdot\delta \vec{r}_{\alpha}=0 \]
其中 $\vec{F}_{\alpha}^{ext}$ 是作用在质点 $\alpha$ 上的外力，$\dot{\vec{p}}_{\alpha}$ 是其动量随时间的变化率。在笛卡尔坐标系下，对于 $n=3N$ 个坐标，可写为标量形式：
\[ \sum_{i=1}^{n}(F_{i}^{ext}-\dot{p}_{i})\delta x_{i}=0 \]

\subsection{从达朗贝尔原理推导拉格朗日方程}
1. \textbf{坐标变换}：引入广义坐标 $q_j$ ($j=1, \dots, n-k$)，虚位移为 $\delta x_i = \sum_j \frac{\partial x_i}{\partial q_j} \delta q_j$。\
2. \textbf{广义力项}：
   将虚功表示为广义坐标的形式：
   \[ \sum_i F_i^{ext} \delta x_i = \sum_{i,j} \left( F_i^{ext} \frac{\partial x_i}{\partial q_j} \right) \delta q_j = \sum_j Q_j \delta q_j \]
   其中 $Q_j$ 定义为广义力。
3. \textbf{惯性项变换}：
   考虑 $\sum_i \dot{p}_i \delta x_i$ 项。假设质量为常数，利用以下两个关键恒等式：
   \begin{itemize}
       \item 速度变换：$\frac{\partial x_i}{\partial q_j} = \frac{\partial \dot{x}_i}{\partial \dot{q}_j}$
       \item 算符交换：$\frac{d}{dt} \frac{\partial x_i}{\partial q_j} = \frac{\partial \dot{x}_i}{\partial q_j}$
   \end{itemize}
   代入动能定义 $T = \sum_i \frac{1}{2}m_i \dot{x}_i^2$，通过分部积分的思想可以证明：
   \[ \sum_i \dot{p}_i \frac{\partial x_i}{\partial q_j} = \frac{d}{dt} \frac{\partial T}{\partial \dot{q}_j} - \frac{\partial T}{\partial q_j} \]
4. \textbf{合并结果}：
   由于 $\delta q_j$ 彼此独立，达朗贝尔原理要求每一项系数必须为零：
   \[ \frac{d}{dt}\frac{\partial T}{\partial\dot{q}_{j}}-\frac{\partial T}{\partial q_{j}}=Q_{j} \]
   若系统是保守的（即存在势能 $V$ 使得 $Q_j = -\partial V / \partial q_j$），且 $V$ 与广义速度无关，定义 $L=T-V$，则导出标准形式：
   \[ \boxed{\frac{d}{dt}\frac{\partial L}{\partial\dot{q}_{j}}-\frac{\partial L}{\partial q_{j}}=0} \]

\section*{哈密顿原理 (Hamilton's Principle) 详解}

\subsection{原理陈述}
哈密顿原理指出，对于一个动力学系统，其真实运动轨迹 $q(t)$ 使得作用量泛函：
\[ I = \int_{t_1}^{t_2} L(q_i, \dot{q}_i, t) dt \]
取平稳值，即其变分为零：
\[ \delta I = \delta \int_{t_1}^{t_2} L dt = 0 \]

\subsection{数学推导过程}

1. \textbf{引入路径变分}：
   假设 $q_i(t)$ 是满足极值条件的真实路径。我们引入一条邻近的虚拟路径 $Q_i(t)$：
   \[ Q_i(t) = q_i(t) + \epsilon \eta_i(t) \]
   其中 $\epsilon$ 是无穷小参数，$\eta_i(t)$ 是任意可微函数。
   由于所有路径必须经过相同的始点和终点，因此要求变分在端点为零：
   \[ \delta q_i(t_1) = \delta q_i(t_2) = 0 \quad \text{或} \quad \eta_i(t_1) = \eta_i(t_2) = 0 \]

2. \textbf{展开作用量的变分}：
   计算 $I$ 的变分 $\delta I$：
   \[ \delta I = \int_{t_1}^{t_2} \sum_i \left( \frac{\partial L}{\partial q_i} \delta q_i + \frac{\partial L}{\partial \dot{q}_i} \delta \dot{q}_i \right) dt \]
   注意到 $\delta \dot{q}_i = \frac{d}{dt}(\delta q_i)$。

3. \textbf{分部积分 (Integration by Parts)}：
   对含有速度变分 $\delta \dot{q}_i$ 的第二项进行分部积分：
   \[ \int_{t_1}^{t_2} \frac{\partial L}{\partial \dot{q}_i} \frac{d}{dt}(\delta q_i) dt = \left[ \frac{\partial L}{\partial \dot{q}_i} \delta q_i \right]_{t_1}^{t_2} - \int_{t_1}^{t_2} \left[ \frac{d}{dt} \left( \frac{\partial L}{\partial \dot{q}_i} \right) \right] \delta q_i dt \]

4. \textbf{应用边界条件}：
   由于端点固定（$\delta q_i(t_1) = \delta q_i(t_2) = 0$），边界项 $\left[ \dots \right]_{t_1}^{t_2}$ 消失。代回原式得：
   \[ \delta I = \int_{t_1}^{t_2} \sum_i \left( \frac{\partial L}{\partial q_i} - \frac{d}{dt} \frac{\partial L}{\partial \dot{q}_i} \right) \delta q_i dt = 0 \]

5. \textbf{导出拉格朗日方程}：
   由于 $\delta q_i$ 是彼此独立的任意变分，根据变分法基本引理，被积函数括号内的各项必须分别恒等于零，从而导出：
   \[ \boxed{\frac{d}{dt} \left( \frac{\partial L}{\partial \dot{q}_i} \right) - \frac{\partial L}{\partial q_i} = 0} \]

\subsection{物理意义与应用}
\begin{itemize}
    \item \textbf{统一性}：哈密顿原理不仅适用于力学，还可推广至电磁学、相对论和量子力学。
    \item \textbf{坐标无关性}：由于原理表述不依赖于特定的坐标系，因此它直接证明了拉格朗日方程在点变换下的协变性。
    \item \textbf{约束处理}：在处理带约束的系统时，哈密顿原理可以通过拉格朗日乘子法方便地引入约束力，如圆盘滚动问题中的摩擦力求解。
\end{itemize}
\section{约束力与拉格朗日乘子法 ($\lambda$-method)}

\subsection{理论表述}
当系统存在 $k$ 个完整约束 $f_j(q_1, \dots, q_n, t) = 0$ 时，坐标不再独立。引入拉格朗日乘子 $\lambda_j$ 后，方程修改为：
\[ \frac{d}{dt} \frac{\partial L}{\partial \dot{q}_i} - \frac{\partial L}{\partial q_i} = \sum_{j=1}^k \lambda_j \frac{\partial f_j}{\partial q_i} \]
右侧项 $\sum \lambda_j \frac{\partial f_j}{\partial q_i}$ 即为系统受到的\textbf{广义约束力}。

\subsection{例题：圆盘沿斜面滚动}
一个质量为 $M$、半径为 $R$ 的圆盘在倾角为 $\alpha$ 的斜面上纯滚动。
\begin{enumerate}
    \item \textbf{拉格朗日量}：设质心位移为 $s$，转角为 $\theta$。
    \[ L = \frac{1}{2}M\dot{s}^2 + \frac{1}{2}I\dot{\theta}^2 + Mgs\sin\alpha \]
    其中 $I = \frac{1}{2}MR^2$。
    \item \textbf{约束条件}：纯滚动要求 $f(s,\theta) = s - R\theta = 0$。
    \item \textbf{运动方程}：
    \begin{itemize}
        \item $M\ddot{s} - Mg\sin\alpha = \lambda \frac{\partial f}{\partial s} = \lambda$
        \item $\frac{1}{2}MR^2\ddot{\theta} = \lambda \frac{\partial f}{\partial \theta} = -R\lambda$
    \end{itemize}
    \item \textbf{求解}：结合 $\ddot{s} = R\ddot{\theta}$，解得加速度 $\ddot{s} = \frac{2}{3}g\sin\alpha$，约束力（摩擦力） $\lambda = -\frac{1}{3}Mg\sin\alpha$。
\end{enumerate}



\section{虚功原理 (Virtual Work)}
系统处于平衡态的充要条件是所有外力在满足约束的虚位移下所做的总功为零：
\[ \delta W = \sum_j Q_j \delta q_j = 0 \]
对于保守系统，这等价于势能 $V$ 的变分为零：$\delta V = 0$。

\section{例题：悬链线 (Catenary)}
求固定两端、长度为 $l$ 的均匀链条在重力作用下的形状。
\begin{itemize}
    \item \textbf{泛函}：最小化总势能 $V = \int \rho g y \sqrt{1+y'^2} dx$。
    \item \textbf{约束}：长度 $\int \sqrt{1+y'^2} dx = l$。
    \item \textbf{变分处理}：引入乘子 $\lambda$，构造辅助函数 $f = (y+\lambda)\sqrt{1+y'^2}$。
    \item \textbf{推导}：由于 $f$ 不显含 $x$，利用 Beltrami 恒等式 $f - y' \frac{\partial f}{\partial y'} = C$，解得轨迹为双曲余弦函数：
    \[ y = a \cosh\left(\frac{x-x_0}{a}\right) - \lambda \]
\end{itemize}

\section{拉格朗日方程的协变性 (Invariance)}
拉格朗日方程在坐标变换（点变换）下形式不变。即若从 $q_i$ 变换到 $s_i(q, t)$，新方程仍为：
\[ \frac{d}{dt}\frac{\partial L}{\partial\dot{s}_{i}}-\frac{\partial L}{\partial s_{i}}=0 \]
这使得我们可以根据问题的几何对称性灵活选择坐标系（如球坐标或柱坐标）。

\section{勒让德变换与哈密顿量的定义}

\subsection{勒让德变换 (Legendre Transformation) 的数学原理}
哈密顿力学的出发点是将拉格朗日量 $L(q, \dot{q}, t)$ 转换为以动量 $p$ 为自变量的函数。
设函数 $f(u)$，定义新变量 $v = \frac{\partial f}{\partial u}$。勒让德变换构造新函数 $g(v)$：
\[ g = uv - f \]
其全微分满足：
\[ dg = v du + u dv - df = v du + u dv - v du = u dv \]
由此可见 $g$ 是 $v$ 的函数，且 $u = \frac{\partial g}{\partial v}$。

\subsection{从拉格朗日量到哈密顿量}
在力学中，我们对速度 $\dot{q}_i$ 进行变换，定义广义动量：
\[ p_i = \frac{\partial L}{\partial \dot{q}_i} \]
哈密顿量 $H(q, p, t)$ 定义为：
\[ H(q_i, p_i, t) = \sum_{i=1}^{n} p_i \dot{q}_i - L(q_i, \dot{q}_i, t) \]

\section{哈密顿正则方程 (Hamilton's Canonical Equations)}

\subsection{方程的推导}
考虑 $H$ 的全微分：
\[ dH = \sum_i \left( \frac{\partial H}{\partial q_i} dq_i + \frac{\partial H}{\partial p_i} dp_i \right) + \frac{\partial H}{\partial t} dt \]
同时从定义式 $H = \sum p_i \dot{q}_i - L$ 求微分：
\[ dH = \sum_i (p_i d\dot{q}_i + \dot{q}_i dp_i) - \left( \sum_i \frac{\partial L}{\partial q_i} dq_i + \sum_i \frac{\partial L}{\partial \dot{q}_i} d\dot{q}_i + \frac{\partial L}{\partial t} dt \right) \]
由于 $p_i = \frac{\partial L}{\partial \dot{q}_i}$，消去 $\sum p_i d\dot{q}_i$ 项：
\[ dH = \sum_i \dot{q}_i dp_i - \sum_i \dot{p}_i dq_i - \frac{\partial L}{\partial t} dt \]
（注：利用了拉格朗日方程 $\frac{\partial L}{\partial q_i} = \frac{d}{dt}\frac{\partial L}{\partial \dot{q}_i} = \dot{p}_i$）。
对比系数得到 \textbf{哈密顿正则方程}：
\[ \boxed{\dot{q}_i = \frac{\partial H}{\partial p_i}, \quad \dot{p}_i = -\frac{\partial H}{\partial q_i}} \]
以及时间导数关系：
\[ \frac{\partial H}{\partial t} = -\frac{\partial L}{\partial t} \]

\subsection{相空间 (Phase Space)}
哈密顿力学在 $2n$ 维相空间中描述系统，状态点由 $(q, p)$ 唯一确定。正则方程描述了状态点在相空间中的演化。

\section{哈密顿力学的应用实例}

\subsection{一维简谐振子}
拉格朗日量：$L = \frac{1}{2}m\dot{x}^2 - \frac{1}{2}kx^2$。
1. 计算动量：$p = \frac{\partial L}{\partial \dot{x}} = m\dot{x} \Rightarrow \dot{x} = \frac{p}{m}$。
2. 构造哈密顿量：
\[ H = p\dot{x} - L = p(\frac{p}{m}) - \left( \frac{1}{2}m(\frac{p}{m})^2 - \frac{1}{2}kx^2 \right) = \frac{p^2}{2m} + \frac{1}{2}kx^2 \]
3. 正则方程：
\[ \dot{x} = \frac{\partial H}{\partial p} = \frac{p}{m}, \quad \dot{p} = -\frac{\partial H}{\partial x} = -kx \]

\subsection{中心力场中的运动}
在极坐标 $(r, \theta)$ 下：
\[ L = \frac{1}{2}m(\dot{r}^2 + r^2\dot{\theta}^2) - V(r) \]
1. 广义动量：$p_r = m\dot{r}, \quad p_\theta = mr^2\dot{\theta}$。
2. 哈密顿量：
\[ H = \frac{p_r^2}{2m} + \frac{p_\theta^2}{2mr^2} + V(r) \]
3. 运动方程：
\[ \dot{r} = \frac{p_r}{m}, \quad \dot{p}_r = \frac{p_\theta^2}{mr^3} - \frac{\partial V}{\partial r} \]
\[ \dot{\theta} = \frac{p_\theta}{mr^2}, \quad \dot{p}_\theta = 0 \quad (\text{角动量守恒}) \]

\section{循环坐标与对称性}
若 $H$ 不显含某个坐标 $q_k$（即 $\frac{\partial H}{\partial q_k} = 0$），则称 $q_k$ 为循环坐标（Cyclic Coordinate）。
由正则方程：
\[ \dot{p}_k = -\frac{\partial H}{\partial q_k} = 0 \Rightarrow p_k = \text{Constant} \]
这表明每一个循环坐标都对应一个动量守恒定律。

\section{泊松括号 (Poisson Brackets)}

\subsection{定义与性质}
对于相空间函数 $f(q, p, t)$ 和 $g(q, p, t)$，泊松括号定义为：
\[ \{f, g\} = \sum_{i=1}^n \left( \frac{\partial f}{\partial q_i} \frac{\partial g}{\partial p_i} - \frac{\partial f}{\partial p_i} \frac{\partial g}{\partial q_i} \right) \]
性质：
1. 反称性：$\{f, g\} = -\{g, f\}$。
2. 线性：$\{af+bg, h\} = a\{f, h\} + b\{g, h\}$。
3. 莱布尼茨法则：$\{fg, h\} = f\{g, h\} + g\{f, h\}$。
4. 雅可比恒等式：$\{f, \{g, h\}\} + \{g, \{h, f\}\} + \{h, \{f, g\}\} = 0$。

\subsection{运动方程的括号形式}
任意物理量 $f$ 的时间演化：
\[ \frac{df}{dt} = \frac{\partial f}{\partial t} + \{f, H\} \]
如果 $f$ 不显含时间且 $\{f, H\} = 0$，则 $f$ 是运动常数。

\section{正则变换 (Canonical Transformations)}

\subsection{动机与定义}
寻找坐标变换 $(q, p) \to (Q, P)$，使得在新坐标下方程仍保持正则形式：
\[ \dot{Q}_i = \frac{\partial K}{\partial P_i}, \quad \dot{P}_i = -\frac{\partial K}{\partial Q_i} \]
其中 $K(Q, P, t)$ 是新的哈密顿量。

\subsection{生成函数 (Generating Functions)}
根据哈密顿原理 $\delta \int (p\dot{q} - H) dt = 0$ 和 $\delta \int (P\dot{Q} - K) dt = 0$，新旧变量通过生成函数 $F$ 联系：
\[ \sum p_i \dot{q}_i - H = \sum P_i \dot{Q}_i - K + \frac{dF}{dt} \]
常见的四类生成函数：
1. $F_1(q, Q, t)$: $p_i = \frac{\partial F_1}{\partial q_i}, \quad P_i = -\frac{\partial F_1}{\partial Q_i}$。
2. $F_2(q, P, t)$: $p_i = \frac{\partial F_2}{\partial q_i}, \quad Q_i = \frac{\partial F_2}{\partial P_i}, \quad K = H + \frac{\partial F_2}{\partial t}$。
3. $F_3(p, Q, t)$: $q_i = -\frac{\partial F_3}{\partial p_i}, \quad P_i = -\frac{\partial F_3}{\partial Q_i}$。
4. $F_4(p, P, t)$: $q_i = -\frac{\partial F_4}{\partial p_i}, \quad Q_i = \frac{\partial F_4}{\partial P_i}$。

\section{哈密顿-雅可比理论 (Hamilton-Jacobi Theory)}

\subsection{H-J 方程的导出}
目标是寻找一个正则变换，使得新哈密顿量 $K \equiv 0$。则新坐标 $Q, P$ 均为常数。
令生成函数为 $S(q, P, t)$（即 $F_2$ 型），代入 $K = H + \frac{\partial S}{\partial t} = 0$：
\[ H\left( q_1, \dots, q_n, \frac{\partial S}{\partial q_1}, \dots, \frac{\partial S}{\partial q_n}, t \right) + \frac{\partial S}{\partial t} = 0 \]
这就是 \textbf{哈密顿-雅可比方程}。$S$ 被称为哈密顿主函数（Principal Function）。

\subsection{哈密顿特征函数}
若 $H$ 不显含时间，令 $S(q, P, t) = W(q, P) - Et$，其中 $E$ 是能量。代入 H-J 方程：
\[ H\left( q, \frac{\partial W}{\partial q} \right) = E \]
$W$ 称为哈密顿特征函数。

\section{刘维尔定理 (Liouville's Theorem)}
定理指出：相空间中给定体积元内的状态点密度在随系统演化时保持不变。
数学表述：
\[ \frac{d\rho}{dt} = \frac{\partial \rho}{\partial t} + \{\rho, H\} = 0 \]
这保证了统计力学中系综分布的稳定性。

\section{总结与对比}
\begin{center}
\begin{tabular}{|l|l|l|}
\hline
特性 & 拉格朗日力学 & 哈密顿力学 \\
\hline
变量 & 广义坐标 $q, \dot{q}$ & 广义坐标与动量 $q, p$ \\
方程 & $n$ 个二阶微分方程 & $2n$ 个一阶微分方程 \\
空间 & $n$ 维构形空间 & $2n$ 维相空间 \\
数学核心 & 变分法 & 勒让德变换 / 辛几何 \\
\hline
\end{tabular}
\end{center}

\end{document}